\originalchapter{磁力、磁矩与材料}\label{ch:magneticforce}
\originalsection{洛伦兹力与回旋运动}
同时存在电场与磁场时, 低速点电荷运动满足$\vct F=q(\ee+\vct v\times\bb)$. 磁力与瞬时速度垂直, 所以$\vct F_B\cdot\vct v=0$, 不能单独改变粒子动能. 这不妨碍它改变运动方向. 电场可以作功.

在均匀磁场中, 分解$\vct v=\vct v_\parallel+\vct v_\perp$. 沿场速度不变, 垂直部分作匀速圆周运动. 向心条件为$mv_\perp^2/R=|q|v_\perp B$, 故
\[
R=\frac{mv_\perp}{|q|B},\quad \omega_c=\frac{|q|B}{m},\quad T_c=\frac{2\pi m}{|q|B}.
\]
两部分合成螺旋线, 每圈轴向前进$v_\parallel T_c$. 正负电荷回旋方向相反. 在高速区应使用相对论动量, 这里公式限于非相对论模型. 质谱仪以电势差给予动能, 再用磁场曲率求质量电荷比; $R$单独不能区分质量与电荷.

\originalsection{电流所受磁力}
载流导线中多个运动电荷的磁力相加, 可写$\dd\vct F=I\dd\vct\ell\times\bb$. 直导线长$\ell$处于均匀场, 总力$\vct F=I\vct\ell\times\bb$. 有限曲线应逐元积分. 两平行长导线距离$d$, 一根在另一根处产生$B=\mu_0I_1/(2\pi d)$, 所以另一根的单位长度受力为$\mu_0I_1I_2/(2\pi d)$. 同向电流相吸, 反向相斥.

看似磁力不作功而磁力又能推动导线, 原因在于导线内载流子、晶格和电源共同交换能量. 移动导线产生感应电动势, 电源必须额外做功或机械能减少, 能量最终来自完整系统的电场与电源, 不是磁力对单个电荷凭空作功.

\originalsection{磁偶极子与悬浮}
平面线圈面积$A$, 电流$I$, 其磁矩$\vct m=IA\uv n$, 单位A$\cdot$m$^2$, $\uv n$依电流右手方向. 均匀外场中每对相对边的力合为零, 但产生力偶. 对矩形线圈, 一对边力大小$I\ell B$, 力臂投影为另一边长乘$\sin\theta$, 相乘得$\tau=IAB\sin\theta$. 因而$\vct\tau=\vct m\times\bb$, 旋转势能为$U=-\vct m\cdot\bb$.

场不均匀时, 对足够小且磁矩固定的线圈, $\vct F=\nabla(\vct m\cdot\bb)$. 这来自微小位移中势能的变化. 在外部无电流的磁静态区域, $\nabla\times\bb=0$, 可等价写为$(\vct m\cdot\nabla)\bb$. 在任意有外电流区域直接替换这两形式可能遗漏条件. 磁矩若随场改变, 也要包括材料响应而不能直接沿用固定磁矩的势能.

材料内的原子磁矩可导致顺磁、抗磁和铁磁响应. 顺磁的诱导磁矩倾向沿外场, 抗磁倾向反向, 铁磁能保留磁化并有磁滞, 所以简单常数磁导率并非普遍规律. 线性各向同性弱磁材料可写$\vct M=\chi_m\vct H$, $\bb=\mu_0(\vct H+\vct M)$. 本册不把材料的微观量子解释列为前置.

通电线圈的非均匀场可对磁体产生向上力, 但“合力为零”只保证平衡, 是否稳定还要看微小偏离后的恢复力. 导电环靠近运动磁体时还会感生电流而产生阻碍磁通变化的力; 这种悬浮应在感应规律下分析, 不能假设普通静止金属自动拥有永久磁矩.


条形磁铁的外部场线从N端走向S端, 内部闭合回到N端. 切开磁铁得到两个小偶极, 不会分别得到孤立N和S. 地磁在粗略尺度上也可用倾斜偶极描述, 当地方向与地理坐标还受实际场结构影响; 指南针只是沿当地水平磁场分量取向.

固定磁矩$m_z$在轴场$B_z(z)$中受$F_z=m_z\dd B_z/\dd z$. 对齐场时趋向较强场, 反向时趋向较弱场, 但反向固定磁矩还可能受转动力矩而翻转. 所以用两个普通磁体抵消竖直重力并不足以保证三维稳定.

小而弱、线性、各向同性的抗磁样品, 忽略退磁修正时可用$\vct m=\alpha_m\bb$, $\alpha_m\simeq\chi_m\mathcal V_s/\mu_0<0$. 建立诱导磁矩时响应从零线性增长, 有效机械磁势为$U_m=-\alpha_mB^2/2$, 不是把$\vct m(B)$直接代入固定磁矩式$-\vct m\cdot\bb$. 也可把局部位移中的力写为$F_i=\sum_{j=x,y,z}m_j\partial_iB_j=\alpha_m\partial_i(B^2)/2$, 其中$i$为$x,y,z$之一, 从而积分回上述机械势. 因而抗磁体趋向较弱场. 加重力后$U_{\rm tot}=Mgz-\alpha_mB^2/2$, 只有它在各允许方向均有局部极小, 才可能稳定悬浮. 此处没有声称任何线圈几何都能实现这种极小; 普通导体的涡流悬浮还属于下一章的动态感应机制.
